\papersection{Data}{sec:data}

\subsection{Bloomberg macroeconomic events}\label{sec:bbg}

The event calendar is the Bloomberg U.S.\ economic release calendar from 2010 to 2026. It
contains 179 release lines---a line is one published statistic, such as ``CPI MoM'' or ``Change
in Nonfarm Payrolls''---and 27,137 release rows, each with the release date and time in Eastern
Time, the actual value and the median of the Bloomberg survey of forecasters. Most price and
labor-market data are released at 08:30, survey indicators such as the ISM indices at 10:00,
and FOMC decisions at 14:00.

The DTW model uses the calendar for its timing, not its values: a release is an instant $T$ at
which information arrives, and the model compares the market's reaction at $T$ with its reactions
at earlier releases. The actual values and the consensus are not used.

Many lines are released together as parts of one report. The CPI report alone produces eight
lines at the same second, and the Employment Situation report twelve. Trading each line
separately would trade the same minute several times, so \ref{sec:universe} groups the
lines into release families before any other step. In the family construction, 163 lines
collapse into 67 families.

\subsection{Futures data (S\&P 500, Nasdaq-100, Dow, 10-year and 2-year notes)}\label{sec:futures}

We use one-minute OHLCV bars from Databento (CME Globex, \texttt{GLBX.MDP3}) from July 2010 to
June 2026. Table~\ref{tab:contracts} summarizes the contracts. \ES, \NQ\ and \YM\ are continuous
front-month series rolled on the calendar; \ZN\ is rolled on volume. For \ES, the first years of
the vendor file had many incomplete sessions, and we re-pulled 2010--2015 so that the whole
sample has complete sessions. \YM\ (\$5 per index point, tick one point) is the only other
U.S.\ equity-index future on the same feed with one-minute history from 2010. Its file has 5.2
million one-minute bars; cleaning keeps 3,754 sessions (84--95\% of each year's sessions), and
4.8\% of the kept minutes are filled with the last price, 8--15\% before 2016, like early \NQ,
and under 2\% from 2018.

\begin{table}[!tb]
  \centering\footnotesize
  \caption{Contracts.}\label{tab:contracts}
  \begin{threeparttable}
  \begin{tabular}{lccccc}
    \toprule
    & \ES & \NQ & \YM & \ZN & \ZT \\
    \midrule
    Underlying & S\&P 500 & Nasdaq-100 & Dow Jones Ind. & 10-year note & 2-year note \\
    Continuous series & calendar roll & calendar roll & calendar roll & volume roll & --- \\
    Tick size & 0.25 points & 0.25 points & 1 point & 1/64 point & --- \\
    Value of one point & \$50 & \$20 & \$5 & \$1,000 & --- \\
    Value of one tick & \$12.50 & \$5.00 & \$5.00 & \$15.625 & --- \\
    One tick, in bps of contract value & $\approx 0.5$ & $\approx 0.1$--$0.3$ & $\approx 0.25$--$1$ & $\approx 1.4$ & --- \\
    Role in this paper & traded & traded & traded & cost-gated & dropped \\
    \bottomrule
  \end{tabular}
  \begin{tablenotes}\footnotesize
    \item Notes: Databento \texttt{GLBX.MDP3}, one-minute bars, July 2010 to June 2026.
    \ZT\ was pulled but dropped: before 2019, fewer than 7\% of its release windows have a
    trade in every minute, and a check against tick-level trades confirmed that the gaps are
    minutes without trading, not missing data.
  \end{tablenotes}
  \end{threeparttable}
\end{table}

\paragraph{Sessions.} A trading session runs from 18:00 to 17:00 Eastern Time and is dated by
the day on which it closes. We drop exchange holidays and keep a session only if it has at least
85\% of the asset's median number of bars per weekday session in that year; the threshold is
relative because \ZN\ and early \NQ\ and \YM\ trade fewer minutes than \ES. Minutes without a trade inside
a kept session are filled with the last price. Bars are labelled by their start: the bar
labelled 08:30 covers 08:30:00--08:30:59.

\paragraph{Rolls.} A path or trade that straddles a change of contract would include the
calendar-spread jump. We drop any game in which the contract changes between the first minute
of its path and its exit, and buy-and-hold benchmarks skip the price jump at each roll.

\begin{figure}[!tb]
  \centering
  \includegraphics[width=\linewidth]{fig01_markets_overview_2x2}
  \caption{\ES, \NQ, \YM\ and \ZN, 2010--2026. Roll-adjusted daily prices rebased to 100. Grey: the
  2019--2026 period; red: the 2022 inflation shock.}
  \label{fig:markets}
\end{figure}

Figure~\ref{fig:markets} shows the equity and Treasury futures over the sample, including \YM.

\subsubsection{Why we chose to look at these}\label{sec:why_assets}

All five contracts are among the most liquid futures in the world, trade nearly around the clock
on CME Globex, and react to the same U.S.\ releases, so one calendar and one session definition
serve them all. They differ in how they should react. A hot inflation print raises the expected
path of the policy rate and lowers the present value of future earnings, more so for the
long-duration technology stocks that dominate \NQ; the Dow behind \YM\ holds fewer technology and
more industrial and financial stocks than the S\&P 500, so it tests whether what works in \ES\
carries to a different mix of the same market. The Treasury notes respond to the same news through
yields: the CPI surprise explains about 30\% of the five-minute reaction of \ZN, as it does for
\ES. What separates the markets for a strategy that trades often is the tick: one tick is about
0.5 bps of contract value in \ES, 0.1--0.3 bps in \NQ\ and 0.25--1 bp in \YM, but about 1.4 bps in
\ZN, large next to its 30-minute moves. \ES\ and \NQ\ are the core of the book; \YM\ was added after
their results were known, as an out-of-sample market; \ZN\ is a candidate that a cost gate admits
or excludes (\ref{sec:gating}).

\subsubsection{Why we dropped the 2-year}\label{sec:why_zt}

The 2-year note future (\ZT) is the most direct read on the expected policy rate and would be the
natural market for rate news. But a path-based method needs a price in every minute of the path,
and before 2019 fewer than 7\% of \ZT's release windows have a trade in every minute. A fresh
pull compared against tick-level trades matched exactly, so the gaps are minutes in which \ZT\ did
not trade, not missing data. With a library that sparse before 2019, the analogues cannot be
built, so we dropped \ZT\ at the data stage.

\subsection{Data for DTW: the event panel}\label{sec:dtwdata}

The DTW model does not use the bars directly. It uses a panel of \emph{games} built from them, one
row for each asset, release family, release time and window (\ref{sec:universe}, \ref{sec:five}).
Each row carries three objects, all taken from the cleaned one-minute closes of
\ref{sec:futures}:
\begin{itemize}\itemsep0pt
  \item \textbf{The path:} the 31 closes ending at the entry minute, standardized within the
  window (equation~\eqref{eq:path}). For the first window the path ends at $T+1$, so it holds the
  last half hour before the release and the first minute of the reaction.
  \item \textbf{The target:} the log return from the close of the entry minute to the close 30
  minutes later, the return that the trade earns.
  \item \textbf{The labels:} the family, the clock slot, the window, the contract and whether the
  family passes the impact screen in that year (the \emph{admitted} flag).
\end{itemize}
A game is kept only if every minute of its path and of its holding window is inside a kept session,
the release minute and the minute after it have a price, and the same contract is traded from the
first minute of the path to the exit. A missing minute is never interpolated; the game is dropped.
The release times come from the Bloomberg calendar (\ref{sec:bbg}), after the 163 release lines
are grouped into 67 families.

Every game, admitted or not, enters the library of candidate analogues; only admitted games are
traded. The library is therefore much larger than the traded set: on \ES\ there are about 57,000
games in 2010--2026 but 3,420 admitted ones (Table~\ref{tab:panel}). The first three years,
2010--2012, are a burn-in period: they fill the library and give the impact screen its 30 earlier
releases, but no game in them is scored. The DTW features are computed for admitted games only,
from strictly earlier games in the library, so the panel can be built once and then split by year
into the training, validation and test samples of \ref{sec:whattest}. The panel for the four
assets has 223,583 rows.

\begin{table}[!tb]
  \centering\footnotesize
  \caption{The DTW event panel, 2010--2026.}\label{tab:panel}
  \begin{threeparttable}
  \begin{tabular}{lrrrr}
    \toprule
    & \ES & \NQ & \YM & \ZN \\
    \midrule
    Clean one-minute bars (millions) & 5.39 & 5.26 & 5.20 & 5.39 \\
    Games, all 67 families & 57,146 & 55,317 & 54,455 & 56,665 \\
    Families admitted in some year & 8 & 8 & 5 & 12 \\
    Admitted games with DTW features & 3,420 & 2,735 & 2,390 & 5,420 \\
    \quad of which scored in training, 2013--2017 & 820 & 525 & 640 & 1,085 \\
    \quad of which scored in validation, 2018--2020 & 840 & 665 & 460 & 1,130 \\
    \quad of which in the test years, 2021--2026 & 1,760 & 1,545 & 1,290 & 3,090 \\
    Trades in the test years, 2021--2026 (after the gate) & 1,202 & 1,141 & 870 & 1,945 \\
    Median analogue dispersion $\sigma_e$ (bps) & 17.8 & 24.1 & 17.0 & 6.1 \\
    \bottomrule
  \end{tabular}
  \begin{tablenotes}\footnotesize
    \item Notes: A game is one (asset, family, release, window). All games form the library of
    candidate analogues; admitted games are those whose family passes the impact screen in that
    year, and only they are traded. Trades in the test years are admitted games that pass the gate
    of equation~\eqref{eq:gate}, after games at the same minute are merged.
  \end{tablenotes}
  \end{threeparttable}
\end{table}

The panel is small where it matters. Training scores 820 \ES\ games and validation 840, and the
signal we look for is a correlation of a few hundredths; this is why the setting is fixed in
advance and checked once (\ref{sec:robust}). The panel is also uneven across assets: the impact
screen admits more families on \ZN, because the Treasury market reacts to more of the calendar, and
\ZN's analogues disagree by a third as much as \ES's in bps, in line with its smaller moves. \YM\ has the
fewest admitted families: CPI, payrolls, the ISM Manufacturing and Services indices and the FOMC
Minutes, the core of the \ES\ list without Core PCE, Retail Sales and Wholesale Inventories.
